% Source: https://artofproblemsolving.com/community/c3h2655301_about_something_wrong_with_the_hilbert_hotel
% Stable identity: aops-topic-2655301-q01
A student claims that the real numbers in the interval $[0,1)$ can be matched one-to-one with the positive integers, contrary to Cantor's theorem. The argument has two parts.
\begin{enumerate}
\item[(i)] ``Cantor's diagonal argument assumes that a real number has countably many decimal digits, but a real number has uncountably many digits.''
\item[(ii)] ``Write $x=0.d_1d_2d_3\dots$ in decimal and cut the string of digits at every digit $9$. Each piece contains no $9$, so it can be read as a base-$9$ numeral; record its value and also its number of leading zeros (so that pieces such as $01$ and $1$ are not confused). This turns $x$ into a list $a_1,a_2,\dots,a_m$ of nonnegative integers. By the Fundamental Theorem of Arithmetic the list can be recovered from the positive integer $2^{a_1+1}3^{a_2+1}5^{a_3+1}\cdots p_m^{a_m+1}$, where $p_m$ is the $m$th prime. Different numbers $x$ give different lists, hence different positive integers.''
\end{enumerate}
Explain what is wrong with each part.
